% =====================================================================
%  FICHE 11 — THEME 2 — Exercices FACILE — Oxydant, réducteur, demi-équations
% =====================================================================
\documentclass[11pt,a4paper]{article}
\usepackage[utf8]{inputenc}
\usepackage[T1]{fontenc}
\usepackage{lmodern}
\usepackage[french]{babel}
\usepackage{amsmath,amssymb}
\usepackage[version=4]{mhchem}
\usepackage{siunitx}
\sisetup{output-decimal-marker={,},per-mode=symbol}
\usepackage{xcolor}
\usepackage{array}
\usepackage{booktabs}
\usepackage{enumitem}
\usepackage[most]{tcolorbox}
\usepackage[a4paper,margin=2.1cm,headheight=26pt,headsep=10pt]{geometry}
\usepackage{fancyhdr}
\usepackage[hidelinks]{hyperref}

\definecolor{primary}{RGB}{146,95,160}
\definecolor{primarydark}{RGB}{92,52,104}
\newcommand{\NO}{\ensuremath{\mathrm{N.O.}}}

\newcounter{exo}
\newtcolorbox{exo}[1][]{enhanced,breakable,boxrule=0.9pt,arc=2pt,colback=primary!4,
  colframe=primary,fonttitle=\bfseries,coltitle=white,left=3mm,right=3mm,top=2.2mm,bottom=2.2mm,
  title={Exercice~\refstepcounter{exo}\theexo\ifx\relax#1\relax\relax\else\ \textmd{--- \itshape #1}\fi}}

\pagestyle{fancy}\fancyhf{}
\fancyhead[L]{\small\textcolor{primarydark}{\textbf{Fiche 11 · Thème 2} — Oxydant, réducteur, demi-équations}}
\fancyhead[R]{\small\textcolor{primarydark}{\textsc{Facile}}}
\fancyfoot[C]{\small\thepage}
\renewcommand{\headrulewidth}{0.8pt}
\renewcommand{\headrule}{\hbox to\headwidth{\color{primary}\leaders\hrule height \headrulewidth\hfill}}

\begin{document}
\begin{center}
{\LARGE\bfseries\color{primarydark}Thème 2 — Niveau Facile}\\[2pt]
{\small\itshape Vocabulaire du redox, demi-équations simples, sens du transfert d'électrons.}\\[-4pt]
{\color{primary}\rule{\textwidth}{1pt}}
\end{center}

\begin{exo}[compléter les définitions]
Complète chaque phrase par le mot juste (\emph{augmente/diminue}, \emph{perd/capte},
\emph{oxydant/réducteur}, \emph{oxydée/réduite}) :
\begin{enumerate}[label=(\alph*),leftmargin=2em,topsep=1pt,itemsep=1pt]
  \item Lorsqu'une espèce s'oxyde, son \NO\ \underline{\phantom{augmente}} et elle \underline{\phantom{capte}} des électrons.
  \item L'oxydant est l'espèce qui \underline{\phantom{capte}} des électrons ; c'est celle qui est \underline{\phantom{réduite}}.
  \item Le réducteur est l'espèce qui \underline{\phantom{capte}} des électrons ; c'est celle qui est \underline{\phantom{réduite}}.
  \item Dans une demi-équation, les électrons s'écrivent toujours du côté de l'\underline{\phantom{oxydant}}.
\end{enumerate}
\end{exo}

\begin{exo}[trois couples classiques]
Pour chacun des couples suivants, écris la \textbf{demi-équation} de réduction et indique le nombre
$n$ d'électrons échangés :
\[
\ce{Fe^{3+}}/\ce{Fe^{2+}},\qquad \ce{Cu^{2+}}/\ce{Cu},\qquad \ce{Ag+}/\ce{Ag}.
\]
Précise à chaque fois qui est l'oxydant et qui est le réducteur.
\end{exo}

\begin{exo}[repérer l'oxydation — d'après livre QCM 11]
Parmi les quatre demi-équations suivantes, laquelle (ou lesquelles) représente(nt) une
\textbf{oxydation} ? Justifie par la variation du \NO.
\[
\text{(1) } \ce{Cu(s) -> Cu^{2+} + 2 e^-} \qquad
\text{(2) } \ce{Cu^{2+} + 2 e^- -> Cu(s)}
\]
\[
\text{(3) } \ce{Fe^{2+} -> Fe^{3+} + e^-} \qquad
\text{(4) } \ce{Fe^{3+} + e^- -> Fe^{2+}}
\]
\end{exo}

\begin{exo}[qui est réduit ? — d'après livre QCM 8]
On considère la réaction : $\ce{H2 + S -> H2S}$.
\begin{enumerate}[label=(\alph*),leftmargin=2em,topsep=1pt,itemsep=1pt]
  \item Donne le \NO\ de l'hydrogène et du soufre dans chaque espèce.
  \item Quel élément est \textbf{réduit} ? Lequel est \textbf{oxydé} ?
  \item Quel réactif joue le rôle d'oxydant ?
\end{enumerate}
\end{exo}

\begin{exo}[identifier Ox et Red par le \NO]
Pour chaque couple, identifie la forme oxydée et la forme réduite en comparant les nombres
d'oxydation :
\[
\ce{Sn^{4+}}/\ce{Sn^{2+}},\qquad \ce{Cl2}/\ce{Cl^-},\qquad \ce{MnO4^-}/\ce{Mn^{2+}}.
\]
\end{exo}

\end{document}
